\documentclass[10pt,letterpaper,twocolumn]{article}

\usepackage[utf8]{inputenc}
\usepackage[english]{babel}
\usepackage{amsmath}
\usepackage{amsfonts}
\usepackage{amssymb}
\usepackage{graphicx}
\usepackage[backend=bibtex, sorting=none]{biblatex}
\usepackage[toc,page]{appendix}
\usepackage{listings}
\usepackage{caption}
\usepackage{array}
\usepackage{algorithm}
\usepackage{algorithmic}

% remove those when done
\usepackage[colorinlistoftodos, textwidth=\marginparwidth]{todonotes}
\setlength{\marginparwidth}{2cm} % for the todonotes

\addbibresource{../docs/refs.bib}

\author{\makebox[.9\textwidth]{Guilherme Aramizo Ribeiro}\\PhD Student 
\and Patcharapol Gorgitrattanagul \\Teaching Assistant \and Dr. Jason R Blough\\Instructor}
\title{Dynamic Measurement and Signal Processing \\Lab 6: Order Tracking}

\begin{document}
\maketitle

\begin{abstract}
    Hello! \todo{Some comment}
    
\end{abstract}

\section{Background and Objectives}

	The purpose of this lab is to understand and apply order tracking algorithms in real data. A Fast Fourier Transform (FFT) based and a Time Varying Discrete Fourier Transform (TVDFT) are used to visualize the order of a test bench motor. This section will overview important aspects of rotating machine testing and the two order tracking methods.
	
	As presented in \cite{blough1}, Order Tracking analyses the non-stationary frequency components associated to a rotating machine in operation. It differs from standard Digital Signal Processing (DSP) techniques due to the non-stationary nature of the frequency of interest. In other words:
	
	\begin{equation}
	X(t) = A(k,t)\sin(2 \pi i (k/p) t + \Phi _k)
	\end{equation}	 

	\begin{tabular}{r c m{5 cm}}
	where & $ A $ & is the amplitude \\ 
	 & $ \Phi $ & is the phase angle of order $ k $ \\ 
	 & $ p $ & is the period of primary order \\ 
	 & $ t $ & is time \\ 
	 & $ k $ & is the tracked order \\ 
	\end{tabular} 
	\todo{rotating machine overview}.
	
	Order refers to the signal component whose frequency is a factor of the primary frequency. It can vary in amplitude and frequency over time. The amplitude variation causes errors in its tracking because all the order tracking methods consider a semi-constant amplitude over time.
	
	The FFT order tracking \todo{background of FFT OT}.
	
	\begin{algorithm}
		\begin{algorithmic}[1]
			\STATE $ acc,\ tach \leftarrow sample(blocksize) $
			\STATE Another line
			\STATE Do this
			\STATE Do that $ a = 2 $
		\end{algorithmic}
		\caption{FFT Based Order Tracking}
		\label{algo:fft}
	\end{algorithm}
	
    \begin{equation}
        R = \dfrac{V_{high} - V_{low}}{2^Q}
        \label{eq:resolution}
    \end{equation} \todo{Show FFT equation}
	
	The TVDFT order tracking \todo{background of TVDFT OT}.
	
    \begin{equation}
        R = \dfrac{V_{high} - V_{low}}{2^Q}
        \label{eq:resolution}
    \end{equation} \todo{Show TVDFT equation}

\section{Apparatus}
    The material used in the lab activity is listed bellow:

    \begin{itemize}
		\item Speed adjustable DC motor with an unbalanced mass; \todo{name}
		\item Tachometer; \todo{name}
		\item Accelerometer;  \todo{name}
		\item Signal conditioner; \todo{name}
		\item Accelerometer Hand Calibrator; \todo{name}
		\item Data Acquisition NI cDAQ 9172;
		\item Data Acquisition NI 9234;
		\item Desktop computer with MATLAB, LabVIEW and device drivers installed \ldots
    \end{itemize}


\section{Experimental Procedures}
	A test bench motor with speed adjust is instrument with a tachometer sensor and mounted on a base. The interface between the base and the ground has a layer of foam that allows the platform to oscillate vertically. An 1-axis accelerometer is glued into the center of the platform.
	
	Initially, the accelerometer is calibrated with the \todo{insert calibrator name here}. At $ 149.2\ Hz $ the calibrator oscillates by $ 1 g RMS $. The square rooted auto-power at $ 149.2\ Hz $, $ A_{f = 149.2\ Hz} $, was measured at $ X\ V $\todo{fill value}. The calibration factor, $ c $ was then calculated as
	
	\begin{equation}
	c = \frac{1 g}{A_{f = 149.2\ Hz}} = \frac{X}{X} = X
	\end{equation} \todo{fill value}
	
	The acquisition parameters were selected as to capture $ 32 $ blocks of $ 476 $ samples, measured at $ 1.651\ kHz $, calculate the auto-power with a Flat Top window. The window size is such that the vibrating frequency falls within the center of a frequency bin to avoid leakage. The sampling frequency is higher than the Nyquist rate.
	
	 
	
	On the center of the base, a 1-axis accelerometer 
    The lab activity contains two components: an in-lab experimental procedure and a simulation performed with the MATLAB software. Both activities are described in detail in the lab manual, so this section describes the exceptions found in the activity.
    
    It wasn't possible to decrease the signal generator output amplitude such that it would be mixed in the noise floor. Because the amplitude value is adjusted in discrete steps and the minimal value is well above the noise floor.


\section{Data Summary}
    A sinusoidal signals with frequency of $356\ Hz$ and amplitude of $1.24\ V_{RMS}$ were acquired during the laboratory section. The time history of the experimental and MATLAB generated data was added to the Appendix, Figure \ref{fig:range5}, \ref{fig:range300} and \ref{fig:range075}. The experiment run using a range of $\pm 5\ V$ is presented in the frequency domain in the Figure \ref{fig:spectral5}.
    
    \begin{figure}[h]
        \centering
        \includegraphics[width=0.8\linewidth]{img/sample.png}
        \caption{Overlayed plot of experimental and simulated data in the frequency domain for $\pm 5\ V$ of range and 8, 12 and 24 number of bits.}
        \label{fig:spectral5}
    \end{figure}


\section{Interpretation and Analysis}
    Equation \ref{eq:resolution} demonstrates that a broad input range or low number of bits decrease the resolution. The time history of the simulation data confirms this statement: Figure \ref{fig:range5} and \ref{fig:range300} both show that the samples are closer to the nominal curve as the number of bits increase, and this distance is smaller in the experiment with the smaller range, $\pm 5\ V$. Also, voltage saturation was observed in the Figure \ref{fig:range075}, where the simulated signals saturates on $\pm 0.75\ V$.
    
    Another point from Figure \ref{fig:spectral5} is that the noise floor of the experimental data is higher than the simulated 24 bits counterpart, even thought the acquisition system also has a 24 bits precision. The experimental data also accounts with noises in the environment and in the DAQ electronics. 
    
    The oscilloscope's experimental SNR analysis was not possible because the signal generator couldn't reduce the sine amplitude low enough. But according to their respective data-sheet, the oscilloscope has a 8 bit resolution, while the acquisition board has 24 bits. And both systems have $\pm 5\ V$ as the narrower acquisition range. So according to Equation~\ref{eq:resolution}, the minimum signal amplitude detectable is $19\ mV$ and $298\ \eta V$, respectively. This follows that the NI system has higher SNR.
    
    Finally, in respect to the critical thinking section, for a sine wave with period of $T = 56 \mu s$ and RMS voltage of $A_{RMS} = 0.35\ V$, the desired sampling rate, $F_s$, voltage range, $L$, and number of bits, Q, for the requirements described in the manual are
    \begin{equation}
    F_s = 524288\ Hz \approx 2^{nextpow2(20/T)}
    \end{equation}
    
    \begin{equation}
    L = \pm 1\ V \approx \pm 2\ (A_{RMS}\ \sqrt 2)
    \end{equation}
    
    \begin{equation}
    Q = 8 > log_2 20
    \end{equation}
    
    being $n = nextpow2(x)$ a function that returns the next power of 2 of x, or in other words, an integer n, $2^{n-1} < x \leq 2^n$.
    
\section{Conclusions}
    On this assignment a physical signal was measured thought a data acquisition system in order to observe the effects of quantization. This data was also compared to simulated signals on MATLAB. Finally, the studies of acquisition parameters were explored with an example case.
    
    From the frequency analysis of the discretized sine waves, higher SNR was observed in signals with higher number of bits and higher input voltage band usage. But experimental data has an inherent process noise floor that might be more significant than the quantization error, thus limiting the SNR to a maximum value. Check \cite{einstein}. Testing saving capabilities.

\printbibliography

\onecolumn
\clearpage
\appendix
\appendixpage

\section{Effects of Number of Bits and Voltage Range}
    \begin{figure}[h]
    	\centering
        \includegraphics[width=0.6\linewidth]{img/sample}
        \captionsetup{width=0.6\textwidth}
        \caption{Time history of the experimental data and quantized simulated data for a ADC input range of $\pm 5\ V$}
        \label{fig:range5}
    \end{figure}
    
    \begin{figure}[h]
        \centering
        \includegraphics[width=0.6\linewidth]{img/sample}
        \captionsetup{width=0.6\textwidth}
        \caption{Time history of the experimental data and quantized simulated data for a ADC input range of $\pm 300\ V$}
        \label{fig:range300}
    \end{figure}
    
    \begin{figure}[h]
        \centering
        \includegraphics[width=0.6\linewidth]{img/sample}
        \captionsetup{width=0.6\textwidth}
        \caption{Time history of the experimental data and quantized simulated data for a ADC input range of $\pm 0.75\ V$}
        \label{fig:range075}
    \end{figure}

\clearpage
\section{Main Script}
    \lstinputlisting[language=Matlab]{../src/lab06.m}

\clearpage
\section{Auxiliary DSP Functions}
    \lstinputlisting[language=Matlab]{../src/myDSP.m}



\end{document}
